\documentclass[12pt]{article}
\usepackage[utf8]{inputenc}
\usepackage[T1]{fontenc}
\usepackage{hyperref}
\title{Tangale Language Resources for K-12 Curriculum Development}
\author{Compiled by Claude Sonnet 5 (claude-sonnet-5) for Andr\'as Kornai}
\date{}
\begin{document}
\maketitle
\section{Introduction}
Tangale (ISO 639-3 \texttt{tan}) is a West Chadic (Afro-Asiatic) language of northern Nigeria, spoken mostly in Gombe State. The search results describe nine vowels with contrastive length, two level tones and vowel harmony. Very little material was found; this report says so plainly.
\section{Primary Educational Resources}
No primer, textbook or school programme for Tangale appeared in the search results.
\section{Linguistic and Reference Materials}
\begin{itemize}
\item Jungraithmayr, \emph{A Dictionary of the Tangale Language} (Dietrich Reimer Verlag, 1991), as described in catalogue and Wikipedia references. No online copy was found. \url{https://en.wikipedia.org/wiki/Tangale_language}
\end{itemize}
No orthography guide was found.
\section{Religious and Community Resources}
\begin{itemize}
\item Tangale Bible portions (1920--28) on YouVersion: \url{https://www.bible.com/versions/3241}
\item Joshua Project people-group page, which reports portions from 1920--1929 and a New Testament listing; no complete Bible is listed: \url{https://joshuaproject.net/people_groups/15248}
\end{itemize}
The missionary J.~S.~Hall is reported to have reduced the language to writing.
\section{Digital and Audio Resources}
None found in these searches.
\section{Curriculum Development Context}
Tangale is not known to be taught in schools. Material for a curriculum is very scarce: one scholarly dictionary and early scripture portions. A curriculum would need new primers and an agreed orthography, and these would need to be developed with the community.
\end{document}